% Part: first-order-logic
% Chapter: lindstrom
% Section: introduction

\documentclass[../../../include/open-logic-section]{subfiles}

\begin{document}

\olfileid{mod}{lin}{int}
\section{పరిచయం}

ఈ అధ్యాయంలో, కొన్ని అదనపు నియంత్రణల కింద సంహతత్వ సిద్ధాంతం, అధోముఖ
లొవెన్‌హైమ్--స్కోలెమ్ సిద్ధాంతం వర్తించే తర్కాల్లో మొదటిస్థాయి
తర్కమే గరిష్ఠమైనదని తెలిపే లిండ్‌స్ట్రోమ్ లక్షణీకరణను నిరూపించడమే
మన లక్ష్యం
(\olref[fol][com][com]{thm:compactness},
\olref[fol][com][dls]{thm:downward-ls}). ముందుగా, ఈ సిద్ధాంతం వర్తించే
తర్కాల సాధారణ వర్గానికి మరింత సాధారణమైన లక్షణీకరణ కావాలి. మన
పరిశీలనను \emph{సంబంధాత్మక} భాషలకే పరిమితం చేస్తాం; అంటే,
\tetoken{విధేయ సంకేతాలు}{!!{predicate}s}, వ్యక్తిగత స్థిరసంకేతాలు
మాత్రమే ఉండి, \tetoken{ప్రమేయ సంకేతాలు}{!!{function}s} లేని భాషలకు.

\end{document}
